% Generated draft: requires mathematical and source review.
\documentclass[11pt]{article}
\usepackage[margin=25mm]{geometry}
\usepackage{fontspec}
\setmainfont{Latin Modern Roman}
\usepackage{amsmath,amssymb,mathtools,unicode-math}
\setmathfont{Latin Modern Math}
\usepackage{enumitem,needspace,xurl,hyperref}
\setlength{\parindent}{0pt}
\setlength{\parskip}{5pt}
\setlength{\emergencystretch}{3em}
\newcommand{\N}{\mathbb{N}}
\newcommand{\Z}{\mathbb{Z}}
\newcommand{\Q}{\mathbb{Q}}
\newcommand{\R}{\mathbb{R}}
\newcommand{\C}{\mathbb{C}}
\newcommand{\F}{\mathbb{F}}
\newcommand{\A}{\mathbb{A}}
\newcommand{\PP}{\mathbb{P}}
\newcommand{\E}{\mathbb{E}}
\newcommand{\eps}{\varepsilon}
\newcommand{\dd}{\,\mathrm{d}}
\begin{document}
\section*{UnsolvedMath ID: 11000137}
\textbf{Source problem code:} catalogue entry
\subsection{Definitions and mathematical statement}
Problem 7 asks to devise a matrix model, in the sense of reference [28], for calculating Euler characteristics of the source's ``Arc-complexes.'' Neither [28] nor the definition and indexing of those complexes is supplied.

For terminology, an ordinary surface arc complex has vertices given by appropriate isotopy classes of essential properly embedded arcs, with simplices corresponding to mutually disjoint arc systems; endpoint, marking and boundary conventions determine the complex. The source's capitalized Arc notation might instead involve weighted systems, a quotient, or another specified construction, and no such identification is made from this extract alone. For a finite cell complex the ordinary Euler characteristic is
\[
 \chi(X)=\sum_{j\ge0}(-1)^j c_j,
\]
where \(c_j\) counts \(j\)-cells. Infinite complexes and quotient/orbifold constructions need finiteness, support or weighted-Euler conventions before this formula defines the intended invariant.

The requested matrix model should produce the relevant Euler characteristics through its specified integrals, partition function or formal expansion, but the source gives no matrix space, potential, measure, expansion parameters, surface family or coefficient-extraction rule. Thus one cannot replace the target by the ordinary Euler characteristic of a familiar unquotiented arc complex, or assign a standard Penner/Kontsevich model without verifying that it computes the specified Arc objects.
\subsection{Short English statement}
Construct a matrix-model method for computing the Euler characteristics of the source's Arc-complexes. The complex definition, Euler-characteristic convention and model reference are missing, so no existing Teichmüller matrix model is silently selected as the answer.
\subsection{Sources}
\begin{itemize}
\item Local UnsolvedMath registry, permanent ID \(11000137\), exact source code \texttt{catalogue entry}, Problem 7, PDF page 119. No individual proposer or proposed year is supplied; the title is truncated.
\item Farb, edited volume \emph{Problems on Mapping Class Groups and Related Topics}: \url{https://www.math.uchicago.edu/~farb/papers/mcgbook.pdf}. The volume has 329 explicitly typeset numbered items across chapters, repeated local numbering distinguished by PDF page.
\item Internal reference [28], introduced by ``cf.,'' for the intended matrix-model context; its bibliography entry is absent. The research summary mentions Kontsevich, Penner and other Teichmüller-related matrix models without complete citations and calls this specific request open.
\item Extraction pdf-font-structure, accessed 2026-07-29. Source evidence NEEDS\_REVIEW presents an open problem and requires a current primary-literature audit. Registry status open and classification OPEN-TRIAGE.
\item Dataset attribution: Ulam AI / UnsolvedMath Contributors, CC BY 4.0. Source rights NEEDS\_REVIEW: publicly accessible author/publisher copy, redistribution terms requiring review.
\end{itemize}
\end{document}
